% from ega4-backmatter-contents-fr.tex
\oldpage[IV]{343}
\EGAIVcontentsline{{\texorpdfstring{\protect\hyperlink{ega4.c04.subsec.20-3}{20.3}}{20.3}}}{Composition des pseudo-morphismes}{237}
\EGAIVcontentsline{{\texorpdfstring{\protect\hyperlink{ega4.c04.subsec.20-4}{20.4}}{20.4}}}{Propriétés des domaines de définition des fonctions rationnelles}{244}
\EGAIVcontentsline{{\texorpdfstring{\protect\hyperlink{ega4.c04.subsec.20-5}{20.5}}{20.5}}}{Pseudo-morphismes relatifs}{249}
\EGAIVcontentsline{{\texorpdfstring{\protect\hyperlink{ega4.c04.subsec.20-6}{20.6}}{20.6}}}{Fonctions méromorphes relatives}{252}

\EGAIVcontentsline{\S21}{Diviseurs}{255}
\EGAIVcontentsline{{\texorpdfstring{\protect\hyperlink{ega4.c04.subsec.21-1}{21.1}}{21.1}}}{Diviseurs sur un espace annelé}{255}
\EGAIVcontentsline{{\texorpdfstring{\protect\hyperlink{ega4.c04.subsec.21-2}{21.2}}{21.2}}}{Diviseurs et Idéaux fractionnaires inversibles}{258}
\EGAIVcontentsline{{\texorpdfstring{\protect\hyperlink{ega4.c04.subsec.21-3}{21.3}}{21.3}}}{Équivalence linéaire des diviseurs}{263}
\EGAIVcontentsline{{\texorpdfstring{\protect\hyperlink{ega4.c04.subsec.21-4}{21.4}}{21.4}}}{Images réciproques de diviseurs}{265}
\EGAIVcontentsline{{\texorpdfstring{\protect\hyperlink{ega4.c04.subsec.21-5}{21.5}}{21.5}}}{Images directes de diviseurs}{267}
\EGAIVcontentsline{{\texorpdfstring{\protect\hyperlink{ega4.c04.subsec.21-6}{21.6}}{21.6}}}{Cycle 1-codimensionnel associé à un diviseur}{270}
\EGAIVcontentsline{{\texorpdfstring{\protect\hyperlink{ega4.c04.subsec.21-7}{21.7}}{21.7}}}{Interprétation des cycles positifs 1-codimensionnels en termes de sous-préschémas}{277}
\EGAIVcontentsline{{\texorpdfstring{\protect\hyperlink{ega4.c04.subsec.21-8}{21.8}}{21.8}}}{Diviseurs et normalisation}{280}
\EGAIVcontentsline{{\texorpdfstring{\protect\hyperlink{ega4.c04.subsec.21-9}{21.9}}{21.9}}}{Diviseurs sur les préschémas de dimension 1}{284}
\EGAIVcontentsline{{\texorpdfstring{\protect\hyperlink{ega4.c04.subsec.21-10}{21.10}}{21.10}}}{Images réciproques et images directes de cycles 1-codimensionnels}{289}
\EGAIVcontentsline{{\texorpdfstring{\protect\hyperlink{ega4.c04.subsec.21-11}{21.11}}{21.11}}}{Factorialité des anneaux réguliers}{302}
\EGAIVcontentsline{{\texorpdfstring{\protect\hyperlink{ega4.c04.subsec.21-12}{21.12}}{21.12}}}{Le théorème de pureté de Van der Waerden pour l'ensemble de ramification d'un morphisme birationnel}{304}
\EGAIVcontentsline{{\texorpdfstring{\protect\hyperlink{ega4.c04.subsec.21-13}{21.13}}{21.13}}}{Couples parafactoriels. Anneaux locaux parafactoriels}{313}
\EGAIVcontentsline{{\texorpdfstring{\protect\hyperlink{ega4.c04.subsec.21-14}{21.14}}{21.14}}}{Le théorème de Ramanujam--Samuel}{323}
\EGAIVcontentsline{{\texorpdfstring{\protect\hyperlink{ega4.c04.subsec.21-15}{21.15}}{21.15}}}{Diviseurs relatifs}{329}
\EGAIVcontentsline{}{BIBLIOGRAPHIE}{333}
\EGAIVcontentsline{}{INDEX DES NOTATIONS}{334}
\EGAIVcontentsline{}{INDEX TERMINOLOGIQUE}{336}

\endgroup
